\ExamPart{Trắc nghiệm trả lời ngắn}{2}{Thí sinh trả lời từ câu 1 đến câu 4.}

\NQuestion Gọi $\alpha$ là góc giữa hai đường thẳng $d_1:x-y+1=0$ và $d_2:x+2y-4=0$. Tính $\tan\alpha$.
\AnswerLine{$3$}
\SolutionBlock{Ta có $m_1=1$ và $m_2=-\dfrac12$. Vậy
\[
\tan\alpha=\left|\dfrac{m_2-m_1}{1+m_1m_2}\right|=\left|\dfrac{-\dfrac12-1}{1-\dfrac12}\right|=3.
\]}

\NQuestion Từ các chữ số $0,1,2,3,4,5,6$, lập được bao nhiêu số tự nhiên lẻ có $2$ chữ số khác nhau?
\AnswerLine{$15$}
\SolutionBlock{Chữ số hàng đơn vị phải là số lẻ nên có $3$ cách chọn: $1,3,5$.\\
Với mỗi cách chọn đó, chữ số hàng chục phải khác $0$ và khác chữ số hàng đơn vị, nên có $5$ cách.\\
Do đó số cần tìm là $3\cdot5=15$.}

\NQuestion Giải phương trình $\sqrt{x^2-5x+7}=x-1$.
\AnswerLine{$x=2$}
\SolutionBlock{Điều kiện là $x-1\ge0\Rightarrow x\ge1$. Bình phương hai vế:
\[
x^2-5x+7=(x-1)^2=x^2-2x+1.
\]
Suy ra $-5x+7=-2x+1\Rightarrow 3x=6\Rightarrow x=2$. Thử lại thấy thỏa mãn.}

\NQuestion Một mảnh vườn hình chữ nhật có chiều dài hơn chiều rộng $3\,\text{m}$ và diện tích bằng $40\,\text{m}^2$. Tính chu vi mảnh vườn.
\AnswerLine{$26\,\text{m}$}
\SolutionBlock{Gọi chiều rộng là $x\,(x>0)$, khi đó chiều dài là $x+3$. Theo đề bài:
\[
x(x+3)=40\Rightarrow x^2+3x-40=0\Rightarrow (x+8)(x-5)=0.
\]
Suy ra $x=5$ m. Vậy chiều dài là $8$ m và chu vi mảnh vườn là $2(5+8)=26\,\text{m}$.}
